\section{The Assertion Ladder}\label{sec:ladder}

This section states the pattern, describes each rung, names the smaller patterns that realise
the rungs, and shows that structure, behaviour and requirements are three loads carried by the
same ladder rather than three separate techniques.

\subsection{The pattern}

\begin{description}[leftmargin=2.2cm,style=nextline]
  \item[Name] Assertion Ladder (also: Graduated Assertion).
  \item[Intent] Keep a technical domain intact when coding agents contribute, by asserting each
    domain rule at the weakest level that still covers its risk.
  \item[Context] A repository shared by people and coding agents that start every session with
    no memory. Changes arrive faster than anyone can review them closely.
  \item[Problem] How do a domain's rules survive contributors who cannot remember them and
    reviewers who cannot read everything?
  \item[Forces] See Figure~\ref{fig:forces}.
    \begin{itemize}
      \item A rule nobody wrote down is lost at the next session.
      \item Written guidance grows until nobody reads it, and competes with everything else in
        an agent's context.
      \item Every executable check costs effort to write, time in the build, and maintenance
        when the domain changes.
      \item The strongest assertions, those made by the type system, depend on the language.
      \item Too much constraint makes the repository hard to change. Ease of use is a force, not
        an afterthought.
      \item A check that an agent can edit is a check that an agent can weaken.
    \end{itemize}
  \item[Solution] For each domain rule, choose a rung. Climb until the rule's risk is covered,
    then stop. When a rule reaches a checked rung, remove it from the prose so that the
    declaration stays short. Run every check behind one command, and put the checks themselves
    behind a human boundary.
  \item[Consequences] Feedback arrives early and names the fix; the manifest stays short;
    intent can be reviewed separately from code. In exchange there are more tests to maintain,
    the top rung differs between languages, and a rule that climbs further than its risk
    warrants makes change expensive.
\end{description}

\begin{figure*}[t]
  \centering
  \resizebox{\linewidth}{!}{% Two forces acting on the domain, and the ladder standing between them and the domain.
\begin{tikzpicture}[
    font=\small,
    box/.style={draw=#1, line width=0.9pt, fill=#1!7, rounded corners=3pt, align=center, inner sep=6pt},
    arr/.style={-{Stealth[length=2.6mm]}, line width=0.9pt, #1},
  ]
  \node[box=jotblue, minimum width=3.2cm, minimum height=1.6cm] (domain)
    {\textbf{The domain}\\[2pt]\footnotesize rules in one place,\\\footnotesize in its own words};
  \node[box=missed, left=2.6cm of domain, text width=3.6cm] (sprawl)
    {\textbf{Sprawl}\\[2pt]\footnotesize locally plausible changes\\\footnotesize spread one rule\\\footnotesize across many places};
  \node[box=behaviour, right=2.6cm of domain, text width=3.6cm] (erosion)
    {\textbf{Cognitive erosion}\\[2pt]\footnotesize the reviewer cannot read it all;\\\footnotesize the agent cannot remember};
  \node[box=structure, below=1.1cm of domain, text width=8.2cm] (ladder)
    {\textbf{Assertion Ladder}: move each rule out of memory and prose\\into the repository's own checks, as high as its risk warrants};
  \draw[arr=missed] (sprawl) -- node[above, font=\footnotesize]{dilutes} (domain);
  \draw[arr=behaviour] (erosion) -- node[above, font=\footnotesize]{hides} (domain);
  \draw[arr=structure] (ladder) -- node[right, font=\footnotesize]{protects} (domain);
  \draw[arr=structure, dashed] (ladder.west) -| node[pos=0.75, left, font=\footnotesize]{fails fast} (sprawl.south);
  \draw[arr=structure, dashed] (ladder.east) -| node[pos=0.75, right, font=\footnotesize]{shows intent} (erosion.south);
\end{tikzpicture}
}
  \caption{Two forces act on a domain once agents contribute. Sprawl spreads a rule across
  places; cognitive erosion hides the spread from both reviewer and agent. The ladder answers
  both by moving rules out of memory and prose into the repository's own checks.}
  \label{fig:forces}
\end{figure*}

\subsection{The rungs}

\paragraph{\rung{L0 Implied}} The rule exists only in someone's head or in habit. The agent
experiences nothing, and guesses. The guess is usually plausible, which is what makes L0
dangerous: nothing about the result looks wrong. Lean would call this tacit knowledge with no
standard: nothing to follow, nothing to improve, and the cost of relearning paid every session.
A teacher would call it the hidden curriculum.

\paragraph{\rung{L1 Shown}} One complete example embodies the rule: a vertical slice from the
domain to the endpoint, small enough to read whole. Agents, like novices, copy structure before
they understand it \citep{sweller1985worked}, so the slice is the cheapest effective assertion
there is. To a lean eye the slice is standard work shown rather than written: the team's current
best way to add a feature, in a form a contributor can copy. Its weakness is that the learner may
copy the wrong detail, or generalise from an accident of the example.

\paragraph{\rung{L2 Told}} A short declaration states the rule, in a file the agent's tooling
loads at the start of every session \citep{agentsmd}. Prose can say what code cannot: what needs
a person's decision, which of two acceptable choices the team prefers, why a rule exists. It is
the written half of standard work. Its weakness is that it can be forgotten, misread or crowded
out, and that it grows.

\paragraph{\rung{L3 Checked}} An executable rule fails the build when the rule is broken. This
is the rung of architecture tests, contract tests and behavioural scenarios. The agent
experiences a red suite and must work out why. This is the built-in test that
Spear and Bowen found attached to Toyota's standard work \citep{spear1999decoding}: the defect is
found where it is made, by the contributor who made it, before the change moves on to a reviewer,
rather than by a separate inspection further down the stream.

\paragraph{\rung{L4 Explained}} The check's failure names the rule, the offender and the fix. A
bare assertion failure tells the agent \emph{that} it is wrong; an explained failure tells it
\emph{what to do instead}, which is the difference between summative and formative feedback. A
bare failure is a stop. An explained failure is a stop that, like an andon signal, shows where
the problem is and why, and names the countermeasure. It costs a sentence per rule.

\paragraph{\rung{L5 Prevented}} The wrong form does not compile. A sealed hierarchy with an
exhaustive switch, or a map keyed by a closed union type, makes an unmapped error impossible to
write. This is mistake-proofing (poka-yoke) of the preventive kind: the fixture admits the part
only one way, so the mistake is never made, where a check at L3 catches it once it has been. It is
also the rung that depends most on the language.

\subsection{Nine rung patterns}

Each rung is realised by one or more smaller patterns, listed in Table~\ref{tab:rung-patterns}.
Two are not tied to a rung: \pattern{Single Gate}, which runs every check behind one command so
that ``finished'' has one meaning, and \pattern{Human Boundary}, which lists the few changes an
agent may not make alone, including changes to the checks themselves.

\begin{table*}[t]
  \centering\small
  \caption{The rung patterns. Each answers one recurring problem.}
  \label{tab:rung-patterns}
  \begin{tabularx}{\linewidth}{l|l|X|X}
    \hline
    \textbf{Rung} & \textbf{Pattern} & \textbf{Problem} & \textbf{Solution} \\
    \hline\hline
    L1 & \pattern{Worked Slice} & How does a newcomer learn the shape of a feature? & One complete vertical slice, small enough to read whole. \\
    \hline
    L2 & \pattern{Short Manifest} & How is what cannot be tested passed on? & One file of a few hundred words, loaded every session. \\
    \hline
    L3 & \pattern{Executable Rule} & How is a structural rule kept when nobody remembers it? & A test that fails when the rule is broken. \\
    \hline
    L3 & \pattern{Closure Check} & How do two lists stay in step? & Assert that every X has a Y. \\
    \hline
    L3 & \pattern{Honest Double} & How do fast tests stay truthful? & One contract, run against the real implementation and its stand-in. \\
    \hline
    L4 & \pattern{Teaching Failure} & How does a red build lead to the right fix? & The failure names the rule, the offender and the remedy. \\
    \hline
    L5 & \pattern{Unrepresentable Mistake} & How is a mistake stopped before any test runs? & Types that make the wrong form fail to compile. \\
    \hline
    all & \pattern{Single Gate} & When is work finished? & One command runs every rung; green means done. \\
    \hline
    all & \pattern{Human Boundary} & Who may change the ladder itself? & A short list of changes that need a person first. \\
    \hline
  \end{tabularx}
\end{table*}

\pattern{Closure Check} recurs more than any other pattern, so it is worth a paragraph of its
own. An
\emph{every X has a Y} assertion keeps two lists in step without either list knowing about the
other: every domain error has an HTTP status; every repository implementation passes the same
contract; every requirement has a scenario and every scenario cites a requirement. It is the
cheapest way we know to make a growing set safe to extend, and extending sets is most of what an
agent does all day.

\subsection{Three kinds of intent, one ladder}

Architecture tests, behaviour-driven development and structured requirements are usually
presented as separate practices with separate literatures. In a repository shared with agents
they are the same move applied to three kinds of intent
(Table~\ref{tab:intents}). \emph{Structure} says where code may live and what it may touch.
\emph{Behaviour} says what the system does, in the domain's own words. \emph{Requirements} say
why, once, in a form that can be traced. Each is declared, exemplified and enforced, and each
has its own closure check.

\begin{table*}[t]
  \centering\small
  \caption{The same three carriers hold three kinds of intent. Text colour follows the intent
  throughout the paper.}
  \label{tab:intents}
  \begin{tabularx}{\linewidth}{>{\raggedright}p{2.6cm}|X|X|X}
    \hline
    \textbf{Intent} & \textbf{Declared (L2)} & \textbf{Exemplified (L1)} & \textbf{Enforced (L3--L5)} \\
    \hline\hline
    \textcolor{structure}{\textbf{Structure}}: where code may live &
      a placement table in the manifest &
      the move-task slice, domain to route &
      architecture rules; exhaustive error mapping \\
    \hline
    \textcolor{behaviour}{\textbf{Behaviour}}: what the system does &
      one convention line: scenarios cite ids and read given / when / then &
      the work-in-progress scenario &
      the scenarios, and the repository contract \\
    \hline
    \textcolor{requirement}{\textbf{Requirements}}: why &
      \repo{REQUIREMENTS.md}, one EARS sentence per rule &
      ten taskboard rules with ids &
      the requirement closure check \\
    \hline
  \end{tabularx}
\end{table*}

\subsection{EARS as a placement oracle}

The most useful property of EARS in this setting was not one it was designed for. Its keywords
classify a rule, and the class predicts where the rule's code belongs (Table~\ref{tab:ears}).
An \emph{unwanted behaviour} about one entity (``If a requested move is not allowed from the
task's status, then\ldots'') is a guard on that entity that raises a domain error. A
\emph{state-driven} rule that spans many entities (``While the work-in-progress limit is
reached\ldots'') needs a collaborator and belongs in a service. An \emph{event-driven} rule whose
response is to tell someone (``When a task moves to done, the board shall announce\ldots'')
is a domain event published by the service with the effect in a handler. An \emph{optional
feature} (``Where a limit is configured\ldots'') is a field on the one settings object. A
\emph{ubiquitous} rule (``The board shall answer every refused request with a problem
document'') is a cross-cutting invariant: one place, guarded by a check.

\begin{table*}[t]
  \centering\small
  \caption{From EARS pattern to layer, with the taskboard rule that exercises each row.}
  \label{tab:ears}
  \begin{tabularx}{\linewidth}{l|X|X}
    \hline
    \textbf{EARS pattern} & \textbf{Implement as} & \textbf{Taskboard rule} \\
    \hline\hline
    Unwanted, one entity & a guard on the entity raising a domain error & R-FLOW-2: an invalid move is refused \\
    \hline
    State-driven, many entities & a service method with a collaborator & R-WIP-1: the work-in-progress limit \\
    \hline
    Event-driven, tells someone & a domain event; the effect in a handler & R-DONE-1: completion is announced \\
    \hline
    Optional feature & a field on the settings object & the configurable limit \\
    \hline
    Ubiquitous & one place, guarded by a closure or architecture check & R-ERR-1: every refusal is a problem document \\
    \hline
  \end{tabularx}
\end{table*}

The mapping matters because the commonest placement mistake an agent makes is to put a correct
rule in the wrong place, and a behavioural test cannot see that mistake (Section~\ref{sec:evaluation}).
A requirement written in EARS carries its placement with it. We present the mapping as a
hypothesis supported by the ten taskboard rules and by one pilot run, not as a law.

\subsection{Compression and fading}

The solution says to climb only as far as a rule's risk warrants, and to remove a rule from the
prose once a check covers it. Both clauses are about ease of use, which is the force teams
forget first. A manifest that repeats what the tests enforce is waste in the lean sense: it costs
reading time every session and adds nothing the gate does not already guarantee. Standard work
is meant to change as the team learns, and here it changes by subtraction as well as addition:
once a rule is built into a check, the written standard no longer needs to carry it. Instructional
design calls the same move \emph{fading}: scaffolding is withdrawn as the learner no longer needs
it \citep{wood1976scaffolding}. In the receivers the manifest says so of itself: ``anything a
test can check lives in the tests, not here.'' Section~\ref{sec:evaluation} reports what the
added rungs cost in words.
